\section{Conclusion}

In this paper, we introduced \BENCHMARK{}, a comprehensive and challenging benchmark designed to evaluate the long-term memory abilities of chat assistants across five core memory tasks: information extraction, multi-session reasoning, temporal reasoning, knowledge updates, and abstention. Through extensive experiments with both commercial systems and long-context LLMs, we demonstrated the significant challenges posed by \BENCHMARK{}, with current systems exhibiting substantial performance drops. By analyzing key design choices across indexing, retrieval, and reading stages, we proposed effective strategies such as session decomposition, fact-augmented key expansion, and time-aware query expansion, which collectively improved both memory recall and the question answering performance. Our findings highlight the need for more sophisticated memory mechanisms to achieve personalized and reliable conversational AI, and \BENCHMARK{} offers a valuable benchmark to drive future advancements in long-term memory capabilities for chat assistants.


\section*{Reproducibility Statement}

In the paper writing and the subsequent code release, we are committed to enabling other researchers to easily leverage our resources, replicate our results, and build upon our findings. 
We have documented the benchmark construction process in detail in \Cref{section-benchmark-creation} and \Cref{appendix-dataset-construction}, including all the attributes and instructions used to prompt LLMs. We have created and will publicly release the two fixed evaluation datasets, \BENCHMARK\textsubscript{\textsc{S}} and \BENCHMARK\textsubscript{\textsc{M}}. In addition, we will also release the algorithm and source mixture used to create these two datasets, so that future studies could build upon them to create chat histories of any length. Finally, we have meticulously documented all the implementation details of our memory optimization in \Cref{appendix-memory-implementation}, and our code will be released along with the benchmark as well. We believe that these efforts for transparency can help advance the field and foster future research endeavors.

\section*{{Ethics Statement}}

The major artifact released by this work is the \BENCHMARK{} evaluation dataset. To construct the chat history, we utilize ShareGPT\footnote{Accessed via \url{https://github.com/lm-sys/FastChat/tree/main}.}, which has an Apache 2.0 license, and UltraChat\footnote{Accessed via \url{https://github.com/thunlp/UltraChat}.}, which has an MIT license. In the data release, we plan to use an MIT license. Since the questions and the evidence sessions are newly created, we conducted a rigorous human screen process to ensure that the new dataset does not contain personally identifiable information or offensive content. This paper involves human annotators in two places: (1) the dataset construction and filtering (\Cref{section-benchmark-creation}) and (2) the manual analysis of commercial systems (\Cref{section-benchmark-challenging}). The process was mainly conducted by three expert annotators, who are in-house NLP researchers that have more than three years of NLP research experience. All the annotators are properly briefed with the annotation objective, and a discussion among them is conducted to resolve uncertain cases. In total, approximately 400 human hours are spent on the dataset construction and 150 hours on the study of commercial systems. The annotators are paid biweekly or monthly, and their salaries include the working hours dedicated to annotation. The data collection protocol has been determined exempt by the IRB of the author's institution. Finally, \BENCHMARK{} and the studied memory system could induce several societal impacts. Storing and recalling user information could cause personal information leakage, and the lack of a memory ``deletion" operator could harm the system's trustworthiness. It is also possible that the memory mechanism could be leveraged by malicious parties to toxic contents into the datastore, which may cause a jailbreak behavior during inference. To mitigate such behaviors, it is essential to implement moderation mechanisms that monitor the read/write data flow of the memory. We encourage producers of memory-augmented assistant systems to be aware of these potential harmful effects and apply thorough testing and mitigation.

